\section{The Legal Argument}
\label{sec:legal}

\subsection{Reframing: Compact by Design, Not by Accident}

Earlier versions of this analysis framed the ensemble comparison as evidence
that the \ar{} plan is ``typical''---near the ensemble median on compactness
and partisan outcomes.  The real data require a sharper framing.

The \ar{} plan is \emph{not} typical.  It is the most compact valid plan,
or close to it.  For Wisconsin, Georgia, and Pennsylvania, the plan sits at
the 0.1--0.2nd percentile of the cut-fraction distribution.  Fewer than 2 in
1,000 random valid plans are more compact.

This is the correct frame: the algorithm does not try to be average.
It tries to find the minimum edge cut.  The ensemble shows it succeeds.

\subsection{The Maximum-Compactness Certificate}

The \ar{} plan for WI, GA, and PA can be presented to courts with the
following certificate:

\begin{keybox}
\textbf{Maximum-compactness certificate (WI/GA/PA)}

A \gc{} \recom{} ensemble of 1,000 valid redistricting plans was run for
this state.  The \ar{} plan has lower edge-cut fraction than
approximately 99.8--99.9\% of ensemble plans.  Any party claiming the
\ar{} plan is non-compact must identify a more compact alternative.
Such an alternative will be found in fewer than 2 in 1,000 random draws
from the feasible space.

This certificate is reproducible: any party can run an independent \recom{}
chain and verify the percentile.  No expert discretion is required.
\end{keybox}

This certificate is legally stronger than a ``typical plan'' certificate.
A plan at the 50th percentile could be replaced by any of thousands of
alternatives.  A plan at the 0.2nd percentile cannot: it is at the
compactness extremum, and displacing it requires accepting a less compact
map.

\subsection{The Geographic-Inevitability Certificate}

For North Carolina, the appropriate certificate is different:

\begin{keybox}
\textbf{Geographic-inevitability certificate (NC)}

A \gc{} \recom{} ensemble of 1,000 valid redistricting plans was run for
North Carolina.  The \ar{} plan's edge-cut fraction ($0.0973$) is within
$0.3\%$ of the ensemble mean ($0.0970$).  The minimum-cut plan and the
random-sampling median have converged.

This means North Carolina's geography, not algorithm choice, determines
the plan.  Any compact redistricting method applied to NC produces
essentially the same partition.  The plan is geographically inevitable.
\end{keybox}

\subsection{Combining Both Certificates with the B.11 Argument}

The G-track evidence structure is:

\begin{keybox}
\textbf{Primary argument (B.11, B.02)}: The \ar{} plan is constitutionally
required by Article~I~\S2, as the geographic completion of the
Huntington-Hill apportionment.

\textbf{Corroborating evidence (G.1)}: The \ar{} plan is the most compact
valid plan for WI, GA, PA (0.1--0.2nd percentile) and the geographically
inevitable plan for NC (50th percentile, within 0.3\% of ensemble mean).
It is not an outlier on compactness---it is the extremum.
\end{keybox}

This combination is legally decisive.  The constitutional argument establishes
why the plan is required.  The ensemble argument establishes that the plan
could not have been drawn more compactly---the algorithm has found the compact
frontier.

\subsection{The Challenger's Problem}

Under the maximum-compactness frame, any legal challenger faces the following
evidentiary burden:

\begin{enumerate}
  \item They must propose an alternative plan.
  \item That plan must have lower edge-cut fraction than the \ar{} plan.
  \item From the ensemble data, such a plan will exist in at most 2 in 1,000
    random draws from the feasible space.
  \item The challenger must therefore either (a) run an ensemble and find the
    rare more-compact plan, or (b) accept that the \ar{} plan is at the
    compactness extremum.
\end{enumerate}

This is an unusually high evidentiary bar.  It inverts the normal structure
of redistricting litigation, where challengers produce alternative maps.  Here,
the map is at the compactness frontier; producing a more compact alternative
requires exceeding what 99.8\% of valid plans achieve.

\subsection{The \texorpdfstring{\rucho{}}{Rucho} Framework Revisited}

\rucho{} held that no judicially manageable standard for identifying excessive
partisan gerrymandering exists at the federal level.  The majority
acknowledged that maps could be evaluated but held that courts could not
determine how much partisanship is too much.

The \ar{} plan offers a different path: a plan that is not selected for
partisan reasons at all, and that is verifiably the most compact feasible
plan.  The court is not asked to determine whether 7D/7R is the ``right''
outcome for North Carolina.  It is asked to certify that the \ar{} plan was
produced by a process that maximises compactness subject to population balance
and contiguity, and that the ensemble confirms this.

The partisan outcomes that result from the most compact valid plan are not
gerrymandered outcomes.  They are geographic outcomes.
